\section{Background and Related Literature}\label{sec:background}

\subsection*{The task mechanism}

The labour-market consequences of AI run through the task. Whether AI raises or lowers the demand for labour depends on which of the individual activities that make up the job it can complete, complement, or cannot do. In the \citet{acemoglu2019} task-based framework, production allocates a continuum of tasks between labour and capital, with technology acting through two opposing channels. While a displacement effect means capital absorbs tasks workers performed, a reinstatement effect means new tasks restore a comparative advantage for labour. Importantly, productivity gains and labour demand can move apart. ``So-so'' automation that displaces workers while delivering only modest productivity gains, has been shown to depress the labour share even as output rises \citep{acemoglu2020}. 
It is, thus, not enough to say exposure to AI implies harm since a task AI can perform is displaced, whereas a task AI makes a worker faster at is complemented, with each carrying opposite signs for labour demand.
%Exposure is therefore necessary but in no way sufficient for harm, as it were.


Which tasks a technology reaches has changed. Earlier computerisation acted on codifiable tasks, as emphasised in the routine-biased technical change literature \citep{autor2003, goos2007, autor2013}, hollowing out middle-wage clerical and production work while sparing both manual and analytical non-routine activity. Newer advances in generative AI reaches the non-routine cognitive tasks the previous wave could not, so exposure no longer splits down the routine/non-routine divide. The updated fault line is consequently the substitution--complementarity margin, and the literature debates on which dominates. \citet{autor2024} reads AI optimistically, as a tool that could reinstate middle-skill judgement and widen access to expert tasks, whereas \citet{acemoglu2025} is more cautious, projecting an aggregate total-factor-productivity gain of under one per cent over a decade once only currently automatable tasks are counted. The mechanism AI will operate along is less disputed, but the sign and magnitude of the impact is an ongoing empirical point of interest.


\subsection*{Measuring exposure and the reliability problem}

Turning the mechanism into a measure requires judging, task by task, what the technology can do. Successive generations of measures have gone about this differently: engineering bottlenecks read at the occupation-level, treating occupations as indivisible \citep{frey2017}; linking AI applications to distinct human occupational abilities \citep{felten2021}; the textual overlap of patents with job descriptions \citep{webb2019}; and scoring the task itself for its suitability for machine learning \citep{brynjolfsson2018}. Large language models have since changed the capability of the measurement. \citet{eloundou2024} have an LLM rate each O*NET task against a rubric, and show the ratings track expert judgment closely while covering the whole taxonomy in a way which no human exercise can afford. The approach now supports institutional studies. ILO work scores every occupation in the international classification to build a global exposure index \citep{gmyrek2023}, the Tony Blair Institute and the OBR score the whole of O*NET for the UK \citep{obr2025, tbi2024}, and the IMF's complementarity-adjusted index separates displacement from augmentation \citep{pizzinelli2023, cazzaniga2024}.

Each of these measures read exposure off different inputs such that they rank the same occupation differently, and so any estimate built on one inherits that disagreement. Recent evidence has also reassessed this problem by moving from a cross-measure concern into a within-measure one. \citet{yin2026} score an identical task set with four LLMs and find average exposure differs by more than a factor of three across them, with downstream labour-market magnitudes varying roughly 2.4-fold by model. The divergence is systematic since each model carries a persistent tendency to score high or low, such that averaging does not remove it. Still, the acknowledgment of exposure scoring as a measurement exercise and not realised impact, is understated in institutional literature, along with a failure to propagate measurement uncertainty into the estimates built on top of it.

\subsection*{The geography of exposure and Scotland's composition}

Exposure is a property of tasks and tasks bundle into occupations, meaning AI's incidence is uneven across space, though the literature disagrees on how just how uneven at the spatial scales of policy interest. \citet{tbi2024} find limited regional variation in UK exposure, attributing it to the diversified sectoral mix of most regions. Additionally, the OECD finds cross-country exposure ranging from 16--77\% with complementarity rising in regional education \citep{oecd2024}, with European evidence finding automation risk driven by industrial structure and agglomeration \citep{crowley2021}. The disagreement is partly artefactual since coarse administrative geographies average over the functional labour markets where local conditions operate \citep{overman2022}. Therefore, the regional signal appears most clearly in occupational mix and least clearly in coarse exposure averages. Consistent with this, \citet{henseke2025} attribute the recent rise in UK exposure to shifts in occupational composition rather than to within-occupation task change.

Scotland's occupational mix is distinctive in ways that bear directly on exposure.
%\footnote{ See \cref{app:descriptives} for the full document of these differences}. 
Its public sector accounts for around 22\% of employment against roughly 18\% UK-wide, a gap persistent over time, with services accounting for 77.1\% of Scottish GDP against 80.5\% for the UK, with more local-government delivery of services such as social care and correspondingly less weight in the professional, scientific and information services concentrated in London and the South East \citep{scotgov2024, commonslibrary2024}. Financial, insurance and professional services are among the most exposed parts of the UK economy, and caring, elementary and skilled-trade work among the least \citep{obr2025}. These exposed sectors divide on the substitution--complementarity margin, since public administration, health and caring work scores high on complementarity and low on substitutability, whereas professional and financial work is exposed on both. Whether Scotland's lighter concentration in the latter harbours its workforce or merely denies it the productivity gains of complemented work is therefore ambiguous. The fiscal edge is sharper given that public-sector pay in Scotland sits about 5\% above the UK median and has diverged upward since 2019 \citep{cribb2025}, coupled against the backdrop of a chronic investment and productivity gap \citep{williams2025}.

\subsection*{Gaps and open questions}

The synthesis leaves several gaps. Exposure measures are reported as point estimates whose dependence on the scoring model and prompt is acknowledged but rarely quantified, and where a measure is carried into downstream exercises that institutions perform, exposure enters as though it were exact, with no way of telling if the resulting figures are trustworthy or where uncertainty arises. Thus far exposure for the UK is estimated only at coarse geographies and has not been posed as a compositional question for a devolved fiscal authority, so the Scotland--rUK differential, and which part of it survives the measurement fragility above, is absent from the literature. Finally, since no instrument for AI exposure yet matches the commuting-zone design of \citet{acemoglu2020}, credible sub-national causal estimates are unavailable, so what remains tractable for this exercise is the measurement of the exposure differential and an audit of what the forecast implications drawn from it actually rest on. These gaps motivate three questions:

\begin{enumerate}
\item Which exposure estimands are robust to cross-model divergence? Specifically, do regional differentials survive the model-specific bias that destabilises absolute levels?
\item Does occupational composition leave Scotland more or less exposed than the rest of the UK, and where within that composition does the gap arise?
\item What does the exposure differential imply for the Scottish income tax net position, and which of the uncertainties along that chain govern the implication?
\end{enumerate}

The remainder of this dissertation takes up these three questions in turn for the Scotland--rUK case.